
\documentclass{article}
\usepackage{colortbl}
\usepackage{makecell}
\usepackage{multirow}
\usepackage{supertabular}

\begin{document}

\newcounter{utterance}

\centering \large Interaction Transcript for game `matchit\_ascii', experiment `similar\_grid\_1', episode 4 with nemotron{-}3{-}ultra.
\vspace{24pt}

{ \footnotesize  \setcounter{utterance}{1}
\setlength{\tabcolsep}{0pt}
\begin{supertabular}{c@{$\;$}|p{.15\linewidth}@{}p{.15\linewidth}p{.15\linewidth}p{.15\linewidth}p{.15\linewidth}p{.15\linewidth}}
    \# & $\;$A & \multicolumn{4}{c}{Game Master} & $\;\:$B\\
    \hline

    \theutterance \stepcounter{utterance}  
    & & \multicolumn{4}{p{0.6\linewidth}}{
        \cellcolor[rgb]{0.9,0.9,0.9}{
            \makecell[{{p{\linewidth}}}]{
                \texttt{\tiny{[P1$\langle$GM]}}
                \texttt{Jūs piedalāties kopīgā uzminēšanas spēlē.} \\
\texttt{Mērķis ir noskaidrot, vai šis režģis un cits attēls, ko redzu tikai es, ir vienādi.} \\
\texttt{Lūdzu, vispirms aprakstiet savu režģi. Pēc tam es sniegšu savu aprakstu, un mēs varēsim viens otram uzdot jautājumus par režģiem, lai noskaidrotu, vai tie ir vienādi. Tagad sāciet savu īso režģa aprakstu ar "APRAKSTS:", kam seko apraksts. Nepievienojiet neko citu. Šis ir jūsu režģis:} \\
\texttt{X X ▢ X X} \\
\texttt{X ▢ ▢ ▢ X} \\
\texttt{▢ ▢ ▢ ▢ ▢} \\
\texttt{X X ▢ X X} \\
\texttt{X X ▢ X X} \\
            }
        }
    }
    & \\ \\

    \theutterance \stepcounter{utterance}  
    & \multicolumn{4}{p{0.6\linewidth}}{
        \cellcolor[rgb]{0.9,0.9,0.9}{
            \makecell[{{p{\linewidth}}}]{
                \texttt{\tiny{[P1$\rangle$GM]}}
                \texttt{APRAKSTIS: 5x5 režģis. Izpildīti rūtiņas (X) veido plusu zīmi centrā. Konkrēti: pirmā rindiņa — 1., 2., 4., 5. kolonna; otrā rindiņa — 1. un 5. kolonna; trešā rindiņa — tikai 3. kolonna; ceturtā rindiņa — 1., 2., 4., 5. kolonna; piektā rindiņa — 1., 2., 4., 5. kolonna. Visas pārējās rūtiņas ir tukšas.} \\
            }
        }
    }
    & & \\ \\

    \theutterance \stepcounter{utterance}  
    & & & \multicolumn{2}{p{0.3\linewidth}}{
        \cellcolor[rgb]{0.9,0.9,0.9}{
            \makecell[{{p{\linewidth}}}]{
                \texttt{\tiny{[GM$|$GM]}}
                \texttt{abort, first line: APRAKSTIS: 5x5 režģis. Izpildīti rūtiņas (X) veido plusu zīmi centrā. Konkrēti: pirmā rindiņa — 1., 2., 4., 5. kolonna; otrā rindiņa — 1. un 5. kolonna; trešā rindiņa — tikai 3. kolonna; ceturtā rindiņa — 1., 2., 4., 5. kolonna; piektā rindiņa — 1., 2., 4., 5. kolonna. Visas pārējās rūtiņas ir tukšas., but expected APRAKSTS:} \\
            }
        }
    }
    & & \\ \\

\end{supertabular}
}

\end{document}
